\section{Further research questions and exact targets}
\label{sec:research}

The questions below use the new results as starting points. Each asks
for a functional identity, a finite reduction, or a specified algebraic
classification. None is asserted as a proved conclusion of this article.

\subsection{Harmonic powers and generalized harmonic numerators}

\begin{question}[A single generating expression in both orders]
The triangular generator determines every $V_r(u)$ using $2r+1$
steps. Find a useful closed expression for
\[
 \sum_{r\ge0}V_r(u)\frac{t^{2r}}{(2r)!}
\]
in a stated domain, with its singular terms in $t$ separated and its
continuation in $u$ controlled. A satisfactory expression should
recover both the negative odd exponent poles and the uniform finite
Euler-sum span without first expanding all the triangular rows.
The centered Mellin kernel and its even local coefficients give a
concrete route. An identity of formal coefficient series should be
distinguished from a normally convergent two-variable identity.
\end{question}

\begin{question}[Minimal coordinates at a fixed harmonic power]
For each $p$, determine the span of all $F_p(-2r,1/2)$ inside the
finite displayed space $\mathcal E_{p-1}$. Compute rational
coefficient matrices for successive $r$ and reduce them modulo an
explicitly stated set of proved multiple-zeta relations. The first
targets are $p=6,7,8$, where the generating theorem is finite but
individual Euler sums begin to interact nontrivially. The desired
output is a minimal presentation relative to those proved relations,
not an unsupported assertion of arithmetic independence.
\end{question}

\begin{question}[A criterion for regular generalized-harmonic products]
Classify polynomials in $H_n^{(q)}$ whose centered Dirichlet series
are holomorphic at every nonpositive even spectral point. The
asymptotic coefficients in inverse powers of $n+1/2$ give necessary
local tests. Determine which combinations involving even harmonic
orders cancel all obstructing coefficients, and express the answer
by a generating identity rather than an expanding list of examples.
This continues Question~11 of \cite{RepoExactResonant} and should
include both sufficient structural families and a clearly specified
necessity statement.
\end{question}

\begin{question}[Exact cancellation of the first global spectral term]
Find finite combinations of the centered harmonic-power families
whose first spectral derivatives at $s=-2r$ reduce to a finite
Euler-sum and Stieltjes expression. The regular values alone cannot
settle this question: the holomorphic Mellin remainder contributes
after differentiation. A proof should first exhibit an exact
identity among the Mellin kernels that cancels that remainder, then
differentiate the completed family. A first target is a relation
among powers $p\le4$ at $s=0$ or $s=-2$.
\end{question}

\subsection{Shift singularities, monodromy, and primitives}

\begin{question}[The full relative-character singularity pattern]
Extend the explicit scalar amplitudes in Theorem~\ref{hsh:germ}
from words $(s,\{1\}^r)$ to arbitrary independent complex orders
in the relative Hurwitz character of \cite{RepoIndependent}.
At each negative shift, determine a finite set of local power--log
germs and their exact coefficients. Then identify which words have
removable earlier shifts and which exceptional order loci reduce the
functional rank. The source transport is available; the challenge
is organizing the cancellations without hiding them in a recursively
defined remainder.
\end{question}

\begin{question}[A monodromy representation with explicit generators]
The local formula gives the monodromy around each negative integer.
Describe the action of finite products of these loops on the span
of relative-character specializations, including the rational poles
created on nonprincipal sheets. Determine whether a finite-depth
filtration is preserved and compute its associated graded action.
This would give a global continuation statement that retains the
branch information responsible for the exact-radius theorem.
\end{question}

\begin{question}[Global normalized primitives and shift transport]
The local primitive in Section~\ref{hsh:primitives} is explicit in
Lerch functions, with finite polylogarithmic formulas for spectral
jets at zero. Construct global primitives normalized at a positive
shift, and determine how their constants transform under
$a\mapsto a+1$ and under monodromy. A useful next target is a
closed difference equation for the first primitive of $D_1$ and its
first two spectral derivatives, retaining every logarithmic contact
term. Repeated primitives may connect to higher Gamma functions,
but any such identification must fix the same base-point constants.
\end{question}

\subsection{Polygamma moments beyond the two-parameter rows}

\begin{question}[Additional generators that isolate the missing moments]
The exact Bell matrices identify the missing formal coordinates at
weights six and eight. Find an additional Gamma or beta parameter
identity whose mixed derivatives contribute a row outside that span.
The immediate target at weight six is an individual value of
\[
 \sum_{n\ge0}(n+a/2)L_2(n;a)^3
\]
for generic $a$, since the two combinations in
\eqref{bell:sixidentity1}--\eqref{bell:sixidentity2} then separate
the other cubic-weight products. At weight eight, choose candidate
generators by pairing their coefficient rows with the explicit
annihilators in \eqref{bell:eightann} before carrying out numerical
experiments. This gives an exact test of whether a proposed family
adds information.
\end{question}

\begin{question}[A hierarchy of generalized-harmonic tail identities]
Specialize the full summation kernel in
\eqref{bell:kernelbasis} at integral and half-integral $a$.
Classify the resulting multiple-zeta or colored multiple-polylogarithm
identities at successive weights and compare their span with stated
shuffle and stuffle relations. The weight-seven tail-square
evaluation is the first compact example. The target is a finite
family valid at arbitrary weight, together with a transparent map
from Bell rows to convergent harmonic sums, rather than a list of
numerically fitted reductions.
\end{question}

\subsection{Collision constants and higher Stieltjes indices}

\begin{question}[Positive Stieltjes indices with moving gaps]
Extend Theorem~\ref{hc:main} to products of
$\gamma_{m_i}^{(r_i)}(\{x+a_i\})$ with some $m_i>0$.
Their singular pieces contain logarithms, so a meromorphic
partial-fraction decomposition alone is insufficient. Introduce
independent local exponents, integrate before differentiating, and
seek a jointly holomorphic remainder after subtraction of the full
parameter-dependent primitives. The first target is three factors
with indices $(1,0,0)$ along both nested paths of
Theorem~\ref{hc:hierarchy}. The statement should specify the
allowed logarithm branches and every intermediate normalization.
\end{question}

\begin{question}[Composition of renormalized collisions]
The exact correction isolates a holomorphic remainder at the full
collision. Can one define partial-cluster corrections so that merging
two clusters and then a third produces the same normalized result as
the other parenthesization? Raw constant extraction is already
nonassociative in the relevant sense, as the oriented harmonic laws
show. Determine the exact defect and whether a local correction
removes it. Start with four digamma factors and two nested clusters;
the primitive coefficients in \eqref{hc:B} give finite algebraic
data for testing the compatibility equations.
\end{question}

\begin{question}[Global multiple kernels after local collision subtraction]
The present theorem isolates all local gap dependence, leaving the
merged constant $Q_{\boldsymbol r}$. Identify permutation or
argument-derivative combinations for which this global constant
reduces to endpoint data, extending the evaluated example
\eqref{hc:middle-derivative}. For the remaining combinations,
give an explicit finite decomposition into convergent multiple
Tornheim-type kernels. Begin with four factors and small total
argument-derivative order; distinguish a local completion theorem
from an evaluation of the surviving global period.
\end{question}

\subsection{The unresolved Gaussian candidates}

\begin{question}[An exact functional bridge for the frozen $S_6$ and $S_8$]
Retain the current source coefficient vectors and the known formal
nonmembership results \cite{RepoMain,RepoExactResonant}. Seek an
additional iterated-integral, distribution, or functional relation
outside the already excluded finite relation families. Any claimed
proof should supply both an exact rational certificate and the
analytic identity represented by each row, including endpoint
regularization. A new relation search inside an unchanged excluded
span cannot settle the conjecture. The Bell and harmonic identities
here provide exact auxiliary periods, but no bridge from them to
the Gaussian targets has been established.
\end{question}

The most immediate extensions are the fixed-power coordinate matrices,
the first positive-Stieltjes-index nested collision, and a new Gamma
generator selected using the weight-six annihilator. Each begins from
a finite, explicit calculation and asks for an exact functional
identity. The broader monodromy and Gaussian questions require
additional structural input.
